\section{Lean~4 Fundamentals}


Lean~4 is simultaneously a programming language and a proof assistant. Programs are
proofs; types are propositions. This section teaches the minimum Lean~4 needed to write
blockchain proofs.

\subsection{Propositions as Types}

In Lean~4, a proposition is a type in \texttt{Prop}. A proof of proposition \texttt{P}
is a term of type \texttt{P}. If you can construct a term of type \texttt{P}, you have
proved \texttt{P}.

\begin{lstlisting}[language=Lean,caption={Basic propositions in Lean~4}]
-- A proposition: for all natural numbers n, n + 0 = n
theorem add_zero (n : Nat) : n + 0 = n := by
  induction n with
  | zero => rfl
  | succ k ih => simp [Nat.add_succ, ih]
\end{lstlisting}

Line by line:
\begin{itemize}[nosep]
  \item \texttt{theorem} declares a named proof obligation.
  \item \texttt{(n : Nat)} universally quantifies over natural numbers.
  \item \texttt{: n + 0 = n} is the proposition (a type in \texttt{Prop}).
  \item \texttt{:= by} enters tactic mode---an interactive proof language.
  \item \texttt{induction n} performs structural induction on \texttt{n}.
  \item \texttt{rfl} closes a goal where both sides are definitionally equal.
  \item \texttt{simp} applies simplification lemmas from the library.
\end{itemize}

\subsection{Tactic Mode}

Tactics are commands that transform proof goals. The most important tactics for blockchain
proofs:

\begin{table}[H]
\centering
\small
\begin{tabular}{@{}ll@{}}
\toprule
\textbf{Tactic} & \textbf{What it does} \\
\midrule
\texttt{intro h} & Introduce hypothesis \texttt{h} from a $\forall$ or $\to$ goal \\
\texttt{exact e} & Close goal with exact proof term \texttt{e} \\
\texttt{apply f} & Reduce goal to subgoals matching \texttt{f}'s premises \\
\texttt{rfl} & Close goal where both sides are definitionally equal \\
\texttt{simp} & Simplify using registered lemmas \\
\texttt{omega} & Decide linear arithmetic over integers/naturals \\
\texttt{induction x} & Structural induction on \texttt{x} \\
\texttt{cases h} & Case analysis on hypothesis \texttt{h} \\
\texttt{have h : P := ...} & Introduce intermediate lemma \\
\texttt{contradiction} & Close goal from contradictory hypotheses \\
\bottomrule
\end{tabular}
\caption{Core tactics for blockchain formal verification.}
\end{table}

\subsection{Dependent Types}

Lean~4's type system allows types to depend on values. This is essential for expressing
properties like ``a list of exactly $n$ elements'' or ``a set of validators where
$|V| \geq 3f + 1$'':

\begin{lstlisting}[language=Lean,caption={Dependent types for validator sets}]
structure ValidatorSet (n : Nat) where
  validators : Fin n -> PublicKey
  threshold : Nat
  threshold_valid : 3 * threshold + 1 <= n
\end{lstlisting}

The field \texttt{threshold\_valid} is not runtime data---it is a proof that must be
provided when constructing a \texttt{ValidatorSet}. The compiler enforces this at
construction time. There is no way to create an invalid validator set.

%% ============================================================
